\begin{figure}[H]
\centering
\begin{tikzpicture}[>=Stealth,scale=0.94]
  \draw[rounded corners,very thick,gray!65] (-5,-2.35) rectangle (5.6,2.35);
  \draw[->] (-5.4,0) -- (6.35,0) node[right] {eje óptico};
  \draw[very thick] (4.5,-2.0) .. controls (5.15,-1) and (5.15,1) .. (4.5,2.0);
  \node[right,align=left] at (4.8,-1.5) {espejo primario\\parabólico};
  \foreach \y in {-1.35,0,1.35}{
    \draw[->,blue,thick] (-4.8,\y) -- ({4.8-0.15*abs(\y)},\y);
  }
  \draw[red,thick] (4.6,1.35) -- (0.67,0.33);
  \draw[red,thick] (4.8,0) -- (1.0,0);
  \draw[red,thick] (4.6,-1.35) -- (1.33,-0.33);
  \draw[very thick] (1.5,-0.5) -- (0.5,0.5);
  \node[align=center] at (-1.0,0.88) {espejo secundario\\plano};
  \draw[->,thin] (-0.55,0.72) -- (0.66,0.37);
  \draw[->,green!60!black,thick] (0.67,0.33) -- (1.0,2.68);
  \draw[->,green!60!black,thick] (1.0,0) -- (1.0,2.68);
  \draw[->,green!60!black,thick] (1.33,-0.33) -- (1.0,2.68);
  \draw[thick,fill=gray!20] (0.58,2.55) rectangle (1.42,2.9);
  \node[above] at (1.0,2.92) {ocular/detector};
\end{tikzpicture}
\caption{Diseño óptico de un telescopio Newtoniano.}
\end{figure}